\fchapter{Ej 2 Pág 40}

\fsection{\boldit{\textcolor{waveBlue2}{UTILIZA LAS TIC.}}} 
\subsection*{Busca en internet una noticia reciente sobre alguna de las consecuencias 
del cambio climático. Lee la noticia y realiza un comentario con la información 
más relevante.}

\begin{solution}
	Según la \href{https://www.rtve.es/noticias/20251116/diez-fenomenos-extremos-agravado-cambio-climatico-espana/16817452.shtml}{noticia} 
	un reciente \href{https://es.greenpeace.org/es/sala-de-prensa/informes/informe-diez-anos-10-eventos-extremos-como-el-cambio-climatico-golpea-a-espana/}{informe} del CSIC confirma que el cambio climático 
	esta intensificando de manera exponencial los fenómenos extremos en España\dots
	En la ultima década, han sido entre entre \SI{1,3}{\celsius} y \SI{2,2}{\celsius} 
	de lo que habrían sido sin el calentamiento global, los grandes incendios forestales
	han visto aumentar su riesgo meteorológico entre un 15 y un 20\% y la devastadora DANA
	de Valencia en 2024 tuvo unas precipitaciones entre un 10 y un 15\% más intensas 
	debido al cambio climático.
\vspace{1em}
	
	Los 10 eventos analizados han provocado en Españá unas 5000 muertes prematuras y pérdidas
	económicas de 23.000 millones de euros en los ultimos 5 años. El estudio concluye que España
	debido a su posición en el mediterráneo es uno de los países europeos más vulnerables y que 
	estos fenómenos ya no son anecdotas sino la nueva realidad.
\vspace{1em}

	Urge, por tanto, acelerar la reducción de emisiones y reforzar de manera 
	inmediata las medidas de adaptación para evitar que los daños y las víctimas 
	sigan creciendo en los próximos años.
\end{solution}
